% TOP_PROBLEM: {"problemId": "problem.constant-query-locally-decodable-codes-polynomial-length", "canonicalTitle": "Constant-Query Locally Decodable Codes with Polynomial Blocklength", "releaseRank": 258, "primaryDomain": "cryptography_coding_information_optimization", "primaryDomainLabel": "Cryptography, coding, information & optimization", "displayStatus": "open", "releaseStatus": "open"}
% !TeX program = lualatex
% Definition and sources only; this file is not an LLM solution attempt.
\documentclass[11pt]{article}
\usepackage[margin=25mm]{geometry}
\usepackage{fontspec}
\setmainfont{Latin Modern Roman}
\usepackage{amsmath,amssymb,mathtools}
\usepackage{unicode-math}
\setmathfont{Latin Modern Math}
\usepackage{xurl,hyperref}
\hypersetup{colorlinks=true,urlcolor=blue}
\setcounter{secnumdepth}{0}
\setlength{\parindent}{0pt}
\setlength{\parskip}{5pt}
\setlength{\emergencystretch}{3em}
\begin{document}
\section*{\texorpdfstring{Constant-Query Locally Decodable Codes with Polynomial Blocklength}{Constant-Query Locally Decodable Codes with Polynomial Blocklength}}\label{258}
\subsection{Definitions and mathematical statement}
Seek absolute constants $q,C\ge1$ and $\delta>0$ and encodings
$\mathrm{Enc}_n:\{0,1\}^n\to\{0,1\}^{N(n)}$ with $N(n)\le n^C$ up to a fixed constant factor. A randomized local decoder, given an index $i$ and oracle access to a word $y$ with Hamming distance at most $\delta N(n)$ from $\mathrm{Enc}_n(x)$, must make at most $q$ queries and satisfy
\[
 \Pr[\mathrm{Dec}^y(i)=x_i]\ge2/3
\]
for every $x,i,y$ in scope. Constants do not depend on message length. The target is existence or impossibility of such a binary code family; it does not ask merely average recovery of a random bit or decoding from a random rather than adversarial corruption pattern.
\par\smallskip\textit{Formulation and concept references:} \href{https://people.seas.harvard.edu/~salil/pseudorandomness/pseudorandomness-Aug12.pdf}{[S1]}.

\subsection{Short English statement}
Can binary messages be encoded with polynomial overhead so that any chosen message bit remains recoverable from a constant number of probes despite a constant fraction of adversarial errors?
\subsection{Sources}
\begin{itemize}
\item[S1]S. P. Vadhan, Pseudorandomness, Open Problem 7.53 (2012).
\url{https://people.seas.harvard.edu/~salil/pseudorandomness/pseudorandomness-Aug12.pdf}.
\textit{Formulation source.}
\end{itemize}
\end{document}
